%! AP Statistics — Unit 8, Lesson 8.1 — Lesson Plan
\documentclass[10pt]{article}
\usepackage{apstats-article}
\usepackage{apstats-boxes}

\newcommand{\UnitNumberName}{Unit 8: Inference for Categorical Data: Chi-Square \quad}
\newcommand{\LessonNumberName}{Lesson 8.1: Introducing Statistics --- Are My Results Unexpected?}

\begin{document}
\begin{center}
    {\Large\bfseries \CourseName: \SchoolYear \\
    {\small\bfseries \UnitNumberName \LessonNumberName}}
\end{center}

\begin{tcolorbox}[colback=sky, colframe=navy, boxrule=0.9pt, arc=2mm,
    left=2mm, right=2mm, top=1.5mm, bottom=1.5mm]
    \textbf{Primary Objective:} Students will recognize that categorical data can differ
    from what we expect by chance, and they will frame the central question of Unit 8:
    are the differences between observed and expected counts large enough to conclude that
    something real is going on, or could they simply be due to random variation?
    (Big Idea: VAR)
\end{tcolorbox}

\renewcommand{\arraystretch}{1.6}
\begin{skillbox}[Priority Ideas \& Skills]{goldbox}
\begin{minipage}[t]{0.48\linewidth}
    \textbf{AP Skill 1.A --- Selecting Statistical Methods}
    \begin{itemize}
        \item Identify questions suggested by variation between observed and expected
              counts in categorical data (Skill 1.A; VAR-1.J).
        \item Bridge: Units 1--2 described distributions of categorical variables using
              counts and proportions --- Unit 8 asks whether those distributions match
              what we expected.
        \item VAR-1.J.1: Variation between what we find and what we expect to find may
              be random, or it may indicate something real.
    \end{itemize}
\end{minipage}
\hfill
\begin{minipage}[t]{0.48\linewidth}
    \textbf{Key Understandings}
    \begin{itemize}
        \item If a distribution of categorical data perfectly matched the expected
              distribution, there would be no statistical question to ask. Variation
              is what creates the question.
        \item Chi-square tests quantify how surprising the observed variation is under
              the assumption that the expected distribution is correct.
        \item This lesson is a framing lesson (Topic 8.1 is not directly assessed
              on the AP exam) that sets the stage for Lessons 8.2--8.7.
        \item Key shift from Units 6--7: instead of one or two proportions, Unit 8
              compares entire categorical distributions.
    \end{itemize}
\end{minipage}
\end{skillbox}

\begin{skillbox}[{Vocabulary, Concepts, \& Theorems}]{greenbox}
\begin{minipage}[t]{0.48\linewidth}
    \begin{itemize}
        \item \textbf{Observed counts} --- the actual frequencies recorded from data
        \item \textbf{Expected counts} --- the frequencies we would predict if a
              specified distribution were true
        \item \textbf{Chi-square test} --- a family of significance tests that compare
              observed and expected counts to decide whether differences are due to chance
    \end{itemize}
\end{minipage}
\hfill
\begin{minipage}[t]{0.48\linewidth}
    \begin{itemize}
        \item \textbf{Categorical variable (review)} --- a variable whose values are
              labels or categories, not numerical measurements
        \item \textbf{Distribution of a categorical variable} --- the list of all
              possible categories and the proportion (or count) of observations in each
        \item \textbf{Goodness-of-fit (preview)} --- how well observed data match a
              claimed distribution
    \end{itemize}
\end{minipage}
\end{skillbox}

\begin{skillbox}[Activate Prior Knowledge \& Spiral Review]{sky}
\begin{minipage}[t]{0.48\linewidth}
    \textbf{Prior Knowledge Check (3--4 min)}
    \begin{itemize}
        \item Units 1--2: two-way frequency tables, marginal and conditional
              distributions of categorical variables.
        \item Units 1--2: bar charts and segmented bar charts to display
              categorical distributions.
        \item Units 6--7: inference logic --- setting up $H_0$ and $H_a$,
              computing a test statistic, interpreting a $p$-value.
        \item Connect: ``In Units 6--7 we compared one or two proportions.
              What if we need to compare an entire distribution at once?''
    \end{itemize}
\end{minipage}
\hfill
\begin{minipage}[t]{0.48\linewidth}
    \textbf{Connections Forward}
    \begin{itemize}
        \item Chi-square goodness-of-fit test setup $\to$ Lesson 8.2
        \item Carrying out the GOF test $\to$ Lesson 8.3
        \item Chi-square test for homogeneity $\to$ Lessons 8.4--8.5
        \item Chi-square test for independence $\to$ Lesson 8.6
        \item Choosing and applying the right chi-square procedure $\to$ Lesson 8.7
    \end{itemize}
\end{minipage}
\end{skillbox}

\begin{skillbox}[Hook \& Lesson]{sky}
\begin{minipage}[t]{0.48\linewidth}
    \textbf{Hook (5--7 min)}\\
    Show a bag of colored candies. Ask: ``The manufacturer claims 20\% of each of
    5 colors. I counted one bag and got Red: 14, Orange: 8, Yellow: 12, Green: 9,
    Blue: 7. Is this bag unusual, or is this just normal variation?''\\[4pt]
    Ask students: ``How would you decide? What information would you need?''\\[4pt]
    Reveal: We need a way to measure how far the observed counts are from
    what we expected. \textbf{That measurement is the chi-square statistic.}
    The question of Unit 8: \textbf{Are differences between observed and expected
    counts large enough to conclude the distribution is not what we thought?}
\end{minipage}
\hfill
\begin{minipage}[t]{0.48\linewidth}
    \textbf{Lesson Flow (15 min)}
    \begin{enumerate}
        \item Review: categorical variables and expected counts. If 50 candies
              should be 20\% each color, the expected count per color is
              $0.20 \times 50 = 10$.
        \item The gap: observed count minus expected count. A gap of 0 means
              perfect agreement. Larger gaps raise more suspicion.
        \item Preview the chi-square idea: we add up all the gaps (adjusted for
              scale). The bigger the total, the more surprising the data. The
              formula and $p$-value arrive in Lesson 8.2.
        \item Frame the unit arc: GOF test (Lessons 8.2--8.3) $\to$ two-way
              tables: homogeneity (Lessons 8.4--8.5) and independence
              (Lesson 8.6) $\to$ synthesis (Lesson 8.7).
    \end{enumerate}
\end{minipage}
\end{skillbox}

\begin{skillbox}[Explicit Instruction \& Active Monitoring]{sky}
\begin{minipage}[t]{0.48\linewidth}
    \textbf{Direct Instruction (8 min)}
    \begin{itemize}
        \item Show the warm-up coffee-shop data on the board. Walk through
              computing expected counts from a uniform distribution claim
              ($25\%$ each of 4 flavors, $n = 200 \Rightarrow$ expected $= 50$ each).
        \item Compute observed $-$ expected for each category. Ask: ``Which
              category surprises you most? Least?''
        \item Emphasize: a single large gap is noticeable, but chi-square tests
              consider \emph{all} categories simultaneously. The formula
              (coming in Lesson 8.2) weighs each gap relative to the expected count.
    \end{itemize}
\end{minipage}
\hfill
\begin{minipage}[t]{0.48\linewidth}
    \textbf{Active Monitoring}
    \begin{itemize}
        \item Listen for ``The largest count must be the most different from
              expected.'' Redirect: a count of 72 with expected 50 is a gap of 22,
              but a count of 5 with expected 50 would be far more surprising.
        \item Cold-call: ``If I expect 25 of each of 4 outcomes and observe
              25, 24, 26, 25 --- are those differences surprising? Why or why not?''
        \item Check warm-up for fluency with proportions and expected counts
              from Units 1--2.
    \end{itemize}
\end{minipage}
\end{skillbox}

\begin{skillbox}[Group Work \& Differentiation]{redbox}
\begin{minipage}[t]{0.48\linewidth}
    \textbf{Partner Activity (8--10 min)}\\
    Use the activity (preferred study environment survey, $n = 120$). Groups:
    \begin{enumerate}
        \item Compute expected counts (equal distribution among 4 options).
        \item Compute observed $-$ expected for each category.
        \item Decide (without a formal test): do these differences look surprising?
        \item Discuss what additional information would tell you whether the
              differences are statistically significant.
    \end{enumerate}
\end{minipage}
\hfill
\begin{minipage}[t]{0.48\linewidth}
    \textbf{Differentiation}
    \begin{itemize}
        \item \textit{Support (Tier R)}: Provide the expected count calculation
              already completed; students focus on computing and interpreting
              differences.
        \item \textit{On-grade (Tier A)}: Complete all four steps.
        \item \textit{Extension (Tier E)}: Write a paragraph explaining what a
              chi-square test would do with these differences; anticipate Lesson 8.2.
        \item \textit{ELL}: Vocabulary card --- ``observed / observado,''
              ``expected / esperado,'' ``chi-square / ji cuadrado.''
    \end{itemize}
\end{minipage}
\end{skillbox}

\begin{skillbox}[Individual Work \& Assessment]{redbox}
\begin{minipage}[t]{0.48\linewidth}
    \textbf{Exit Ticket (5 min)}
    \begin{enumerate}
        \item A fair spinner has 4 equal sections. It is spun 80 times:
              A=25, B=18, C=22, D=15. How many times would you expect each
              outcome if the spinner is fair?
        \item Which outcome has the largest difference from expected?
              How large is the difference?
        \item Without a formal test, do you think the spinner is unfair?
              What would make you more confident?
    \end{enumerate}
\end{minipage}
\hfill
\begin{minipage}[t]{0.48\linewidth}
    \textbf{AP-Style Check}\\
    \textit{A teacher claims students are equally likely to choose any of 4 topics
    for a project. Of 100 students: 35 chose A, 20 chose B, 28 chose C, 17 chose D.
    What is the expected count for each topic under the teacher's claim?}
    \begin{enumerate}[label=(\Alph*)]
        \item 17, 20, 28, 35
        \item 35, 20, 28, 17
        \item 25, 25, 25, 25
        \item 100, 100, 100, 100
    \end{enumerate}
    Answer: (C)
\end{minipage}
\end{skillbox}

\begin{skillbox}[Reinforcement \& Extension]{goldbox}
\begin{minipage}[t]{0.48\linewidth}
    \textbf{Homework / Practice}
    \begin{itemize}
        \item Problem 1: Candy bag (200 candies, 5 colors, 20\% each claimed) ---
              find expected counts, identify the color with the largest absolute
              difference, and describe what a GOF test would answer.
        \item Problem 2: Fair die (120 rolls) --- compute observed $-$ expected
              for each of 6 faces and decide informally whether the results
              are surprising.
        \item Reflection: describe in your own words what question chi-square
              tests answer about categorical data.
    \end{itemize}
\end{minipage}
\hfill
\begin{minipage}[t]{0.48\linewidth}
    \textbf{Extension / Preview}
    \begin{itemize}
        \item \textit{Extension}: Find a published survey with a claimed categorical
              distribution. Compute expected counts and assess whether the observed
              data appear consistent with the claim.
        \item \textit{Preview}: Lesson 8.2 introduces the chi-square goodness-of-fit
              test --- we will learn the formula for combining all observed vs.\
              expected differences into a single test statistic and $p$-value.
    \end{itemize}
\end{minipage}
\end{skillbox}

\end{document}
